%
% ********** Pagina de Rosto
%

% titlepage gera paginas sem numera�ao
\begin{titlepage}

\begin{center}

\small

% O comando @{} no ambiente tabular x  para criar um novo delimitador
% entre colunas que no a barra vertical | que  normalmente utilizada.
% O delimitador desejado vai entre as chaves. No exemplo, no h nada,
% de modo que o delimitador  vazio. Este recurso est sendo usado para
% eliminar o espao que geralmente existe entre as colunas
\begin{tabularx}{\linewidth}{@{}l@{}C@{}r@{}}
% A figura foi colocada dentro de um parbox para que fique verticalmente
% centralizada em relao ao resto da linha
\parbox[c]{3cm}{\includegraphics[width=\linewidth]{LogoUFRN}} &
\begin{center}
\textsf{\textsc{Universidade Federal do Rio Grande do Norte\\
Centro de Tecnologia\\
Departamento de Engenharia de Computa��o e Automa��o\\
Curso de Engenharia de Computa��o}}
\end{center}
\end{tabularx}


% O vfill  um espao vertical que assume a mxima dimenso possvel
% Os vfill's desta pgina foram utilizados para que o texto ocupe
% toda a folha
\vfill

\LARGE

\textbf{Modelo de estacionamento automatizado utilizando sensoriamento sem fio}

\vfill

\Large

\textbf{Iuri Guerra de Freitas Pereira}

\vfill

\normalsize

Orientador: Prof. Dr. Luiz Eduardo Cunha Leite
% Se no houver co-orientador, comente a pr�xima linha
%\\[2ex] Co-orientador: Prof. Dr. Beltrano Catandura do Amaral

\vfill

\hfill
\parbox{0.5\linewidth}{
% Descomente as opes que se aplicam ao seu caso
\textbf{Monografia}
apresentada � Banca Examinadora do Trabalho de Conclus�o do Curso de Engenharia de Computa��o 
e Automa��o em cumprimento �s exig�ncias legais como requisito parcial � obten��o
do t�tulo de Engenheiro de Computa��o e Automa��o.}

%\textbf{Disserta��o de Mestrado}
%\textbf{Tese de Doutorado}
%apresentada ao Programa de P�s-Gradu��o em Engenharia El�trica da UFRN
%(�rea de concentra��o: Automa��o e Sistemas)
%(�rea de concentra��o: Telecomunica��es)
%como parte dos requisitos para obten��o do t�tulo de
%Mestre em Ci�ncias.}
%Doutor em Ci�ncias.}

\vfill

\large

% Este n�mero de ordem deve ser obtido na coordenao do PPgEE
% Corresponde ao n�mero sequencial da sua tese ou dissertao:
% por exemplo, a 25 tese de doutorado ter o nmero de ordem D25
% Evidentemente, este dado no existe para propostas de tema, caso
% em que a prxima linha deve ser comentada.
%Nmero de ordem PPgEE: M000
Natal, RN, Dezembro de 2010

\end{center}

\end{titlepage}

